\chapter{Amoebiasis}
\index{Amoebiasis}

\section*{Definition and Overview}
Amoebiasis is an infection of the large bowel caused by the single-celled parasite \textit{Entamoeba histolytica}. Many infected people remain symptomless carriers, but the parasite can cause amoebic dysentery, characterised by diarrhoea containing blood and mucus, and in a minority of cases it invades beyond the gut to produce amoebic liver abscess and, rarely, disease of the lungs or brain. The infection is acquired by swallowing the parasite's hardy cysts in contaminated food or water and is most common in regions with poor sanitation. Amoebiasis is eminently treatable with specific antiparasitic drugs, provided the diagnosis is made. It is important to distinguish \textit{E. histolytica} from morphologically identical but harmless commensal species such as \textit{E. dispar}, which do not cause disease.

\section*{Epidemiology}
Amoebiasis occurs worldwide but is most common in tropical and subtropical regions with poor sanitation, including parts of Africa, South Asia, Latin America, and the Pacific. Hundreds of millions of people are infected each year with \textit{E. histolytica} or its harmless look-alike relatives, although only a minority develop invasive disease. In high-income countries the infection is uncommon and is seen mainly in returned travellers, recent migrants, and men who have sex with men. Invasive amoebiasis is more frequent and more severe in young children, older adults, pregnant women, malnourished individuals, and people with weakened immunity. Improved sanitation and safe water are the fundamental public health measures that reduce transmission.

\section*{Aetiology and Pathophysiology}
Amoebiasis is acquired by the faecal-oral route, through swallowing the resistant cysts of \textit{Entamoeba histolytica}:

\begin{itemize}
    \item \textbf{The parasite:} \textit{E. histolytica} has a two-stage life cycle, with motile trophozoites in the gut and hardy, infective cysts passed in faeces.
    \item \textbf{Transmission:} cysts are swallowed in contaminated food or water, or transferred on hands, food, or household surfaces after poor hygiene.
    \item \textbf{Invasion:} in some people the trophozoites adhere to and invade the lining of the large bowel, causing flask-shaped ulcers, inflammation, and bloody diarrhoea.
    \item \textbf{Spread beyond the gut:} invasive trophozoites can reach the portal circulation and travel to the liver, where they form an amoebic abscess, and rarely disseminate to the lungs or brain.
    \item \textbf{Determinants of illness:} many infections remain symptomless; disease is more likely with certain parasite strains, malnutrition, immunosuppression, and other host factors.
\end{itemize}

\section*{Clinical Features}
Many infected people have no symptoms, and illness ranges from mild diarrhoea to severe dysentery:

\begin{itemize}
    \item Diarrhoea with blood and mucus (amoebic dysentery).
    \item Crampy abdominal pain, bloating, and lower abdominal tenderness.
    \item Fever, malaise, and a general feeling of being unwell.
    \item Weight loss and fatigue with more chronic or recurrent infection.
    \item With a liver abscess: fever, pain in the right upper abdomen, and occasionally referred pain in the right shoulder.
    \item Symptoms usually begin one to four weeks after exposure, but they can appear months later, and a liver abscess may develop years after leaving an endemic area.
\end{itemize}

\section*{Differential Diagnosis}
Bloody diarrhoea and right-sided abdominal pain have many other causes:

\begin{itemize}
    \item Bacterial dysentery due to \textit{Shigella}, \textit{Salmonella}, or \textit{Campylobacter}.
    \item \textit{Clostridium difficile} infection, particularly after antibiotics.
    \item Inflammatory bowel disease, including Crohn's disease and ulcerative colitis.
    \item Ischaemic colitis.
    \item Other intestinal parasites, including \textit{Giardia} and cryptosporidium.
    \item Viral gastroenteritis.
    \item Colorectal cancer, which must be considered with chronic or atypical symptoms.
    \item Appendicitis and other causes of right-sided abdominal pain.
\end{itemize}

\section*{When to Seek Medical Care}
See a doctor for diarrhoea that is bloody, severe, lasting more than a few days, or associated with fever, weight loss, or recent travel to a region with poor sanitation.

\textbf{Contact your local emergency services for:}
\begin{itemize}
    \item Severe abdominal pain or a rigid, tender abdomen.
    \item High fever with chills.
    \item Heavy bleeding from the bowel or blood in vomit.
    \item Signs of dehydration, including dizziness, confusion, or very little urine.
    \item An inability to keep fluids down.
\end{itemize}

\section*{Diagnosis and Investigations}
Stool testing is the mainstay of diagnosis:

\begin{itemize}
    \item \textbf{Stool tests:} microscopy can identify cysts or trophozoites, but stool PCR or antigen testing is preferred because it distinguishes \textit{E. histolytica} from harmless look-alike species.
    \item \textbf{Blood tests:} full blood count and inflammatory markers; serological antibody tests are very helpful in diagnosing amoebic liver abscess.
    \item \textbf{Sigmoidoscopy or colonoscopy:} examination of the bowel may show the typical discrete ulcers, with biopsy to confirm invasion.
    \item \textbf{Ultrasound or CT of the abdomen:} to detect a liver abscess and guide management.
    \item \textbf{Travel history:} essential, because the diagnosis is much more likely in returning travellers and migrants.
\end{itemize}

\section*{Management}
Treatment uses two types of drug, one active against parasites invading the bowel wall and one to clear cysts from the lumen:

\begin{itemize}
    \item \textbf{Tissue-active drug:} metronidazole or tinidazole to kill trophozoites invading the bowel and other tissues.
    \item \textbf{Luminal agent:} paromomycin or diloxanide furoate to clear remaining cysts and prevent relapse.
    \item \textbf{Hydration:} oral rehydration solution, with intravenous fluids when dehydration is severe.
    \item \textbf{Pain relief:} simple analgesia as needed.
    \item \textbf{Liver abscess:} most respond to antibiotics alone; percutaneous drainage is reserved for large abscesses or those at risk of rupture.
    \item \textbf{Contact review:} household members may require testing and treatment if they are symptomless carriers.
    \item \textbf{Follow-up:} repeat stool tests to confirm that the infection has cleared.
\end{itemize}

\section*{Complications}
Without adequate treatment, amoebiasis can cause serious complications:

\begin{itemize}
    \item Amoebic liver abscess, which can rupture into the chest or abdomen.
    \item Perforation of the bowel and peritonitis.
    \item Severe dehydration, especially in children.
    \item Rare dissemination to the lungs or brain.
    \item Anaemia and malnutrition with chronic infection.
    \item Narrowing (stricture) of the bowel in long-standing cases.
\end{itemize}

\section*{Prognosis and Outcomes}
With prompt, correct treatment, most people with amoebic dysentery recover fully, and even liver abscesses resolve with antibiotics and, where needed, drainage. Untreated infection can persist for months and carries a risk of serious complications. The disease is rarely fatal when treatment is available, but the risk of death rises substantially if an abscess ruptures or the bowel perforates. Relapse is uncommon when the luminal agent is taken to clear remaining cysts, and repeat stool testing confirms cure.

\section*{Prevention and Self-Care}
Prevention relies on food and water hygiene, because there is no vaccine:

\begin{itemize}
    \item Drink safe water: bottled, boiled, or treated water in areas with poor sanitation.
    \item Eat only food that is well cooked and served hot.
    \item Wash and peel fruit and vegetables, and avoid raw salads in at-risk regions.
    \item Wash hands thoroughly with soap after using the toilet and before eating.
    \item Be cautious with ice and tap water when travelling.
    \item Practise good food hygiene at home if a household member is infected.
    \item Remember that avoiding contaminated food and water is the key protection, as no vaccine exists.
\end{itemize}

\section*{Sources and Further Reading}
The following reputable sources provide reliable, up-to-date information on amoebiasis:

\begin{itemize}
    \item NHS. \textit{Amoebiasis information for patients and travellers.} https://www.nhs.uk/
    \item Centers for Disease Control and Prevention. \textit{Amebiasis information for health professionals.} https://www.cdc.gov/
    \item MSD Manual. \textit{Amebiasis overview.} https://www.msdmanuals.com/
    \item MedlinePlus. \textit{Amebiasis health information.} https://medlineplus.gov/
    \item World Health Organization. \textit{Information on water, sanitation, and parasitic diseases.} https://www.who.int/
    \item Mayo Clinic. \textit{Travelers' health and infectious disease information.} https://www.mayoclinic.org/
\end{itemize}
